\chapter{Controllo Braccio Robotico}
Questo capitolo contiene le principali scelte che sono state fatte per il controllo del braccio robotico. Sì passerà da una prima fase di basso livello, molto tecnologica e dipendente dal tipo di motori che si cuole utilizzare, ad una parte di più "alto livello" che include gli algoritmi con cui il sistema aziona i motori

\section{Azionamento dei motori}
L'azionamento dei motori considerati (MG996R) è abilitato mediante PWM. Più precisamente bisogna considerare un pwm scandito su un periodo di \(20\ ms\), dove per il nostro servomotore vengono colte solo le variazioni che vengono effettuate per pwm agenti tra l' \(1\ ms \) ed i \(2\ ms\).
Dove la variazione viene colta dal motore solo se superiore ad una variazione di \(5\ \mu s\). Sì va anche a considerare il massimo grado di apertura, che dal datasheet ci dice di essere tra -60° e +60°, per cui ha un angolo di apertura di massimo 120°.
Da questi dati, eseguiamo i seguenti calcoli:
\[
angoloMinimo = \frac{5\ \mu s}{1\ ms} * 120 = 0.6°
\]

Tale valore ci interessa dato che per generare il PWM, si farà un utilizzo dei timer per cui, nella prossima sezione, si presenterà il modo tecnico con cui i pwm sono stati implementati

\subsection{Timer per i PWM}
I Timer presenti sulla scheda Nucleo-F401RE hanno una modalità di funzionamento detta proprio PWM. Principalmente quello che si va a fare, settato il periodo (inteso come numero di tick), si può settare un massimo di altri 4 registri, che vanno a confrontarsi con il registro principale di decremento, finchè il tick counter principale non raggiunge il valore effettivo di quello settato nei registri, l'uscita rimane alta, raggiunto, l'uscita sarà bassa fino alla fine del periodo. Pertanto sarà importante settare in maniera opportuna i timer per poterli far funzionare nel modo in cui ci piace.
Prima di partire con la configurazione dei timer, tramite il sistema STM32CubeMX, sì va a visualizzare il clock tree, facendo attenzione al clock che arriva ai timer. Notiamo che il clock che arriva ai timer 2 e 3 (di nostro interesse dato che hanno i canali non in conflitto con i pin che andrà ad utilizzare l'expansion board bluetooth), notiamo che il clock è pari ad 84 MHz. 
Tale valore permette di poter fare il seguente calcolo, un timer ha 2 valori che vanno settati:
\begin{itemize}
	\item \textbf{Prescaler}: Registro che permettere di fare una prima "divisione" della frequenza di clock

	\item \textbf{ARR}: Registro che permette di definire i Tick entro i quali si scandisce un periodo

\end{itemize}

La scelta che è stata fatta parte dalla "sensibilità" del circuito di controllo del servomotore. Data la sua caratteristica, dividendo \(1\ ms\) per \(5\ \mu s\), sì ottiene che il numero di colpi per definire un periodo e che mi permetta di avere la minor sensibilità possibile è 200. Pertanto un periodo sarà scandito da \(200 * 20 = 4000\ Tick\). Da questo dato dobbiamo andare a settare il prescaler in modo che la frequenza passi da \(84\ MHz\) a \(\frac{1}{5\ \mu s} = 200\ kHz\). Eseguendo una divisione tra i clock per determinare il numero di Tick si ottiene:
\(\frac{84\ MHz}{200\ kHz} = 420 Tick\). Da tali calcoli, quindi, andremo a settare sia per il Tim 3 che per il Tim 2 i seguenti valori:
\begin{itemize}
	\item \textbf{Prescalere}: Settato a 420
	\item \textbf{ARR}: Settato a 4000
\end{itemize}

Settati tali valori, la configurazione è pronta, bisogna solo attivare i canali in modo che possano fare il PWM out. Sono stati attivati 2 canali del timer 2 e 3 canali del timer 3. Il pin di uscita è direttamente settato da CubeMX, il che quindi ci permette di capire che è un vincolo. Per questo bisogna anche fare attenzione alla parte di configutrazione del sistema bluetooth

Il settaggio dei registri che pilotano effettivamente il pwm viene effettuata mediante queste linee di codice:

\begin{lstlisting}[language=C, title={Settaggio del resgistro CCR (Pulse) con STMCubeIDE e HAL}]
TIM_OC_InitTypeDef sConfigOC;
sConfigOC.OCMode = TIM_OCMODE_PWM1;
sConfigOC.Pulse = value;
sConfigOC.OCPolarity = TIM_OCPOLARITY_HIGH;
sConfigOC.OCFastMode = TIM_OCFAST_DISABLE;

HAL_TIM_PWM_ConfigChannel(&htimx, &sConfigOC, TIM_CHANNEL_x);
HAL_TIM_PWM_Start(&htimx, TIM_CHANNEL_x);
\end{lstlisting}

Al posto delle x sono stati inseriti i valori effettivi del sistema precendentemente configurato. Da notare che non sì da in input il pin d'uscita ma il timer con il corrispettivo canale. 

La \textbf{Mappatura dei motori} viene effettuata all'interno della libreria "servo.h" che contiene tutte le funzioni e le macro univocamente dedicate al basso livello. Quindi in caso si dovesse cambiare l'hardware del braccio robotico, tale file è da rivedere. Il principale compito di tale libreria è quello di garantire una sicurezza per i motori (evitando, nel caso di richieste follim di andare a rompere i motori), di gestire la loro mappatura rispetto a degli ID più facili. Per gli utilizzatori di tale libreria vi è la funzione "set\_degrees(angolo, id)", che nella sua implementazione va a settare un determinato valore all'interno di quei registri di confronto in base all'ID del motore selezionato. 
Al suo interno effettua una scelta su quale motore andare ad agire e, nel caso non si possa effettuare il movimento richiesto, da come output intero un valore pari ad 1. Che permette di decidere poi, al di fuori del basso livello, come andare poi ad agire

\section{Controllo del braccio}
Il controllo del braccio vero e proprio è stato implementanto cercando di evitare la considerazione dovuta al pwm. Principalmente il funzionamento del controllo viene gestito da una ISR che viene risvegliata ogni 20 ms. Tale ISR va poi a dire effettivamente come il braccio robotico dovrà muoversi.

I controlli progettati sono 2, che abbiamo chiamato:
\begin{itemize}
	\item \textbf{Lineare}: Va semplicemente ad effettuare degli incrementi in base ai dati dell'accellerometro
	\item \textbf{Filtrato}: Va ad effettuare un integrazione in base ai valori che arrivano dall'accellerometro. Rispetto all'approccio lineare permette di poter avere dei movimenti dei motori più morbidi
\end{itemize}

\subsection{ISR periodica}
La ISR implementata è stata fatta mediante una apposita configurazione del timer 4 e del suo prescaler. Precisamente, sappiamo che al timer 4 arrivano \(16\ MHz\) pertanto è importante definire il valore del prescaler e dell'ARR. Ma sopratutto, a differenza del PWM, settare l'accadimento di un evento come se fosse un interrupt. Considerando che non abbiamo un sistema operativo, per eseguire delle operazioni periodiche andiamo ad utilizzare lo scatto di tale timer che va a risvegliare una ISR con una priorità più alta, che va poi ad agire sull'azionamento.
Questo sistema permette di poter utilizzare il massimo della velocità per l'applicazione di un controllo del braccio. (Tale passaggio ha permesso al controllo filtrato di poter essere utilizzato ed applicato).

La ISR è strutturata in 3 parti principali:
\begin{itemize}
	\item \textbf{Bufferizzazione dati}: La libreria di controllo del braccio condivide due variabili con la libreria di ricezione dei dati. Quando vengono ricevuti dei dati, la ISR li bufferizza all'interno di una sua struttura dati. (che utilizzerà come riferimento fino a nuovi dati)
	\item \textbf{Applicazione del controllo}: Si passano i dati all'algoritmo di controllo filtrato che ci andrà ad agire direttamente sui motori, quindi andrà direttamente a controllare il braccio
	\item \textbf{Ripartenza del timer}: Il timer è settato in one Pulse mode, e quindi ripartirà alla fine della ISR effettiva. Tale operazione viene fatta manualmente per evitare che si scatenino delle ISR prima che la precedente sia stata effettivamente servita
\end{itemize}

